% engines: lualatex
% expect-warnings: 0
% expect-errors: 0
\documentclass{article}
\usepackage[margin=1cm]{geometry}
\usepackage[kernel = lua]{../../latex/pgf-go}
\newcommand\case[1]{\par\medskip\noindent\textbf{#1}\par\nopagebreak}
\pagestyle{empty}
\begin{document}

\case{Move 1 captures the marked A1: keep (the triangle stays) / remove (it goes with the stone) / fade and remove}
\begin{goban}[5, captures]
	\stones[w]{A1}\stones[b]{B1}\marks[triangle]{A1, C3}\moves{A2}
\end{goban}
\begin{goban}[5, captures, captured marks = remove]
	\stones[w]{A1}\stones[b]{B1}\marks[triangle]{A1, C3}\moves{A2}
\end{goban}
\begin{goban}[5, captures = fade, captured marks = remove]
	\stones[w]{A1}\stones[b]{B1}\marks[triangle]{A1, C3}\moves{A2}
\end{goban}

\case{Sequence over A1 (captured), D4 (stone), E5 (empty) and C3 (stone): 2, 3 and 4 keep their numbers}
\begin{goban}[5, captures, captured marks = remove]
	\stones[w]{A1,D4,C3}\stones[b]{B1}\marks[sequence]{A1, D4, E5, C3}\moves{A2}
\end{goban}

\case{The square on A1 goes with the stone; the circle on B1 stays (not captured)}
\begin{goban}[5, captures, captured marks = remove]
	\stones[w]{A1}\stones[b]{B1}\marks[w, square]{A1}\marks[b, circle]{B1}\moves{A2}
\end{goban}

\end{document}
